\chapter{Performance and numerical hygiene}
\label{ch:performance}

This chapter records the implementation choices that carry a performance
or numerical-branch consequence: the ordering of the MO integral
transform, the blocked CI solver, the automatic solver selection, the
jit reference-array form, and one repeated contraction in the CI
matrix--vector product.  Each choice is an engineering decision with a
measured anchor; none changes the definition of the energy.  The
measurements were taken on one 64-core node with JAX in double
precision on CPU, with the BLAS thread count fixed at one; wall clock
is quoted as a relative quantity only, while compilation counts and
peak resident memory are load-independent.

\section{MO integral transforms: four sequential contractions}

The two-electron transform is the most expensive single contraction in
the differentiable chain, and its algebraic ordering is part of the
numerical contract.  The implementation is
\texttt{core.nevpt2.impl.stepwise\_mo\_ints}, which performs four
sequential contractions,
\begin{equation}
  h_{ij} = \sum_{ab} C_{ai}\, h_{ab}\, C_{bj},
  \qquad
  (ij|kl) = \sum_{pqrs} C_{pi} C_{qj} C_{rk} C_{sl}\, (pq|rs),
\end{equation}
as one two-index transform followed by three one-index transforms:
\begin{lstlisting}[language=Python]
def stepwise_mo_ints(Cx, h_, e_):
    h1 = jnp.einsum('ai,ab,bj->ij', Cx, h_, Cx)
    z = jnp.einsum('pi,pqrs->iqrs', Cx, e_)
    z = jnp.einsum('qj,iqrs->ijrs', Cx, z)
    z = jnp.einsum('rk,ijrs->ijks', Cx, z)
    eri = jnp.einsum('sl,ijks->ijkl', Cx, z)
    return h1, eri
\end{lstlisting}

Every MO transform in the differentiable chain calls this function,
including the three canonicalization transforms and the state-averaged
driver.  A second copy in the driver delegates to it, so the two
call paths cannot drift apart.

\subsection{Why the one-shot form is not used}

The one-shot form contracts four orbital coefficients against the
four-index integral in a single expression,
\texttt{jnp.einsum('pi,qj,rk,sl,pqrs->ijkl', C, C, C, C, eri)}.  It
is algebraically identical and was the original implementation.  Two
measurements removed it.

First, on a $96^4$ integral array the one-shot form carried an absolute
error around $10^{-6}$ against the sequential transform (the recorded
value is $1.28\times10^{-6}$, with the active block agreeing to
$8.9\times10^{-16}$).  The error is harmless for the energy, which is
insensitive to it, but a linear molecule's generalized Fock matrix
contains exact $\pi$ and $\delta$ degeneracies whose gaps are at the
$10^{-15}$ level.  A perturbation of $10^{-6}$ against a gap of
$10^{-15}$ rotates the eigenvector basis inside the degenerate subspace
completely; the ratio is about $10^9$.  The theta gradient differed by
eight percent between the two contraction orders, and the
Hessian-vector product was NaN for the one-shot form and finite for the
stepwise one.

Second, the sequential form is faster at production size, not slower.
On the physical ThO/ECP array with 96 orbitals ($\max|eri| = 4.78$),
the one-shot evaluation took $5.08$--$5.79$~s and the stepwise
evaluation $0.417$--$0.467$~s, a factor of eleven to thirteen; on a
random array of the same shape both forms finished in about
$50~\mu$s.  The difference is the contraction path XLA selects for the
physical memory layout, and it is a runtime effect, not a
reproducibility precaution.

A release-tree cross-check qualifies the first measurement.  On the
same physical array the current one-shot form agrees with the stepwise
form to $3.33\times10^{-16}$ (rms $5.73\times10^{-18}$), so the
$10^{-6}$ defect was not reproduced in this JAX version and on this
array.  The ordering is retained for both reasons above: the branch
sensitivity of a degenerate subspace is a property of the contraction
order rather than a fixed numerical offset, and the sequential form is
the faster route at the size where the transform dominates.

\subsection{Accuracy on the release tree}

The two orders were compared at fixed geometry and fixed theta, before
and after the substitution.  The differences are at or below
$5.7\times10^{-14}$~Ha for the energy and $2.7\times10^{-14}$ for the
gradient; the non-degenerate H2O case is bit-identical.

\begin{center}
\begin{tabular}{lrr}
\toprule
System and quantity & $\Delta E$ (Ha) & $\Delta g_{\max}$ \\
\midrule
H2O/6-31G CAS(2,2), CASSCF energy and theta gradient & $0$ & $0$ \\
CO/6-31G CAS(4,4), linear, CASSCF energy and theta gradient
  & $5.7\times10^{-14}$ & $2.7\times10^{-14}$ \\
H2O/6-31G CAS(2,2), NEVPT2 energy (eager and jit) & $0$ & --- \\
H2O/6-31G CAS(2,2), NEVPT2 theta gradient & --- & $6.7\times10^{-15}$ \\
\bottomrule
\end{tabular}
\end{center}

The acceptance band for the substitution was $10^{-12}$~Ha on
non-degenerate systems; the largest observed energy difference is
$5.7\times10^{-14}$.

\section{The $S_z$-blocked CI solver}

\texttt{core.casci.cs\_casci\_blocked} diagonalizes the CI Hamiltonian
one spin-projection block at a time.  The partition is exact: the
Hamiltonian is strictly zero between different $S_z$ blocks (the
largest measured off-block element is $0.000\times10^{0}$), so the
global lowest eigenvalue is the lowest eigenvalue over the block
minima.  For H2O/6-31G CAS(2,2) the block structure is one determinant
at $S_z = +6$, four at $S_z = +8$, and one at $S_z = +10$; the global
minimum lies in the four-determinant block, whose smallest in-block
gap is $2.545\times10^{-2}$, and the block energies are separated by
$0.34$~Ha.

The solver exists because the unblocked CI matrix carries a three-fold
degeneracy at $-84.8416418$~Ha for H2O/6-31G CAS(2,2).  The reverse
rule of a dense \texttt{eigh} forms $1/(w_j-w_i)$ for every pair in the
matrix and returns NaN at that degeneracy; inverse iteration defers
the NaN from the first derivative to the second but does not remove it.
Block separation is necessary but not sufficient: the reverse pass
forms the same factor for every pair inside a block whether or not the
pair carries cotangent, so one bit-exact pair anywhere in any block
turns the untouched entries into $0 \times \infty$.  The block
eigenproblem therefore goes through \texttt{core.eigh.eigh\_safe},
whose softened rule freezes the rotation inside a degenerate subspace.

\subsection{Memory and time at two determinant counts}

The two anchor spaces are CAS(6,6) with 924 determinants and CAS(8,8)
with 12870 determinants, both on H2O/6-31G.  The 12870 case used
ncore~=~1.  Peak resident memory was measured under a cgroup limit; the
plain 12870 run required a 48~GB limit and reached 37.83~GB.

\begin{center}
\begin{tabular}{lrll}
\toprule
Route & $E$ (Ha) & Time & Peak RSS \\
\midrule
plain, eager & $-85.21730226017026$ & $27.65$~s & $37.83$~GB \\
blocked, eager & $-85.21730226017024$ & $84.1$--$88.7$~s & $6.78$~GB \\
blocked, jit & same & $41.70$~s ($1.97$ compile $+$ $39.73$ evaluate)
  & $7.78$~GB \\
plain, jit & --- & not run & --- \\
\bottomrule
\end{tabular}
\end{center}

The memory ratio is $5.58$ in eager execution and $4.86$ in the jitted
form.  The time direction is the opposite: blocked eager is $3.2$
times slower than plain eager, and the blocked jit evaluation
($39.73$~s, excluding its $1.97$~s compile) is still $1.44$ times the
plain eager time, although it is $2.1$ times faster than its own eager
form.  The plain jitted form was not run because the same
unblocked CI graph already exceeded 2400~s of compilation at 924
determinants; the 12870-determinant graph is larger.  The blocked
solver's benefit is therefore memory and compilation feasibility, not
single-call time: it builds one block at a time, so the largest boolean
tensor drops from $(12870, 12870, 18)$ to $(4900, 4900, 18)$ at the cost
of per-block dispatch.  The energy difference between the two routes
is $2\times10^{-14}$~Ha.

At 924 determinants the picture reverses: plain takes
$14.7$--$16.9$~s at $1.20$--$1.23$~GB and blocked takes
$27.8$--$29.8$~s at $1.26$--$1.33$~GB.  There is no memory benefit and
the eager solver is $1.8$ times slower.

\section{Automatic solver selection}

\texttt{core.nevpt2.impl.resolve\_ci\_solver} maps
\texttt{ci\_solver='auto'} to one of the two solvers from the
determinant count,
\begin{equation}
  N_{\mathrm{det}} = \binom{2\,n_{\mathrm{act}}}{n_{\mathrm{ele,act}}},
\end{equation}
with the threshold \texttt{BLOCKED\_CI\_MIN\_DET}~$=2000$: at or above
it the blocked solver is selected, below it the plain solver.  An
explicit \texttt{'plain'} or \texttt{'blocked'} always wins, and any
other value raises \texttt{ValueError}.  The threshold is an
engineering choice between the 924 and 12870 anchors, not a measured
optimum; the crossover between the two anchors was not measured.

\begin{lstlisting}[language=Python]
from diffmrc.core.nevpt2.impl import (canonicalize_jax,
                                      resolve_ci_solver)

# Inspect the decision without running a calculation:
solver = resolve_ci_solver('auto', nact, nele_act)   # 924 -> 'plain'

# canonicalize_jax accepts 'auto', 'plain', or 'blocked':
C_can, ci_can, eps = canonicalize_jax(
    C, h1e, eri, ci, ncore=ncore, nact=nact, nele_act=nele_act,
    ci_solver='auto')
\end{lstlisting}

The flag applies to the two internal CI re-solves of the
canonicalization routines (\texttt{canonicalize\_jax} and
\texttt{canonicalize\_sa\_jax}).  The main CI solve of the production
chain remains on the plain solver deliberately: the blocked solver
returns the exact eigenvector, while the plain solver returns an
inverse-iteration vector with a residual of order $10^{-8}$, and
substituting the vector would change the NEVPT2 correlation energy by
more than the $10^{-12}$ acceptance band.  In the state-averaged
routine the flag also has no effect when more than one root is
requested, because there is no multi-root blocked solver and the
state-average branch of \texttt{cs\_casci} is required.

\subsection{What the switch leaves invariant}

Orbitals and orbital energies are bit-identical between the two routes;
the CI vector differs by a global phase and by the plain solver's
residual.  The values below are the phase-aligned differences.

\begin{center}
\begin{tabular}{lrrrr}
\toprule
System & $N_{\mathrm{det}}$ & $\max|\Delta C|$ & $\max|\Delta\epsilon|$
  & CI difference \\
\midrule
H2O/6-31G CAS(2,2) & 6 & $0$ & $0$ & $4.71\times10^{-10}$ \\
CO/6-31G CAS(4,4) & 70 & $0$ & $0$ & $5.39\times10^{-8}$ \\
\bottomrule
\end{tabular}
\end{center}

The identity is the intended semantics: the internal re-solve enters
no returned quantity except the CI vector, while the natural-orbital
rotation and the sub-block Fock diagonalization are built from the
input vector.

\subsection{The small-system penalty the default avoids}

Forcing the blocked solver on a small space is a net cost.  On the
H2O/6-31G CAS(2,2) NEVPT2 chain, the eager correction takes
$23.50$~s with the plain solver and $31.41$~s with the blocked solver,
a $36$ percent increase, and the end-to-end wall clock moves from
$60.3$~s to $74.2$~s.  With \texttt{'auto'} the same chain resolves to
plain and returns to $22.26$~s eager and $62.8$~s wall clock, within
the load band of the pre-change measurement.  The measured anchors
behind the threshold are collected below.

\begin{center}
\small
\begin{tabular}{r l p{0.24\textwidth} p{0.26\textwidth} l}
\toprule
$N_{\mathrm{det}}$ & System & plain & blocked & Decision \\
\midrule
6 & H2O CAS(2,2) & $23.50$~s eager, $60.3$~s wall
  & $31.41$~s eager, $74.2$~s wall & plain \\
70 & CO CAS(4,4) & $9.5$--$16.4$~s & $17.2$--$24.6$~s & plain \\
924 & H2O/CO/N$_2$ CAS(6,6) & $14.7$--$16.9$~s, $1.20$--$1.23$~GB
  & $27.8$--$29.8$~s, $1.26$--$1.33$~GB & plain \\
12870 & H2O CAS(8,8) & $27.65$~s, $37.83$~GB
  & $84.1$--$88.7$~s, $6.78$~GB & blocked \\
\bottomrule
\end{tabular}
\end{center}

\section{Reference arrays as jit arguments}

The theta-only factories compute the reference SCF and the AO integrals
once, at construction, and return a function of the rotation
parameters.  The reference arrays can enter that function in one of two
ways: closed over, or passed as traced arguments.  In the closure form
the arrays become constants of the traced graph, so a jitted wrapper
compiles one executable per geometry.  In the array form the arrays are
runtime arguments, and one executable serves every geometry whose
arrays have the same shape and dtype.

The API exposes both members on the returned callable:
\texttt{E\_theta.component\_args} is the array-argument function and
\texttt{E\_theta.arrays} is the tuple to pass, in order.  The
state-averaged energy factory \texttt{make\_sa\_energy\_fn} carries the
same pair, and the MC-PDFT table passes them through
\texttt{converge\_orbital\_point(..., arrays=E\_sa.arrays)} so that the
orbital search does not compile the AO integrals into its graph.

\begin{lstlisting}[language=Python]
import jax

E_theta = make_nevpt2_theta_fn(mol, ncore, nact, nele, nmo, pairs,
                               scf_builder)
E_jit = jax.jit(E_theta.component_args)   # one executable
E_ref = E_jit(theta, *E_theta.arrays)     # for every same-shaped geometry

# The orbital-point search takes the same pair:
converge_orbital_point(E_theta.component_args, len(pairs),
                       arrays=E_theta.arrays)
\end{lstlisting}

The measured difference on three geometries (the reference scaled by
$-1$ percent, $0$, and $+1$ percent about the center of mass) is
one compile against three, and $2.67$~s against $8.60$~s for the three
calls, measured as interleaved medians.

\begin{center}
\begin{tabular}{lrr}
\toprule
Form & jit objects / compilations & Total for three geometries \\
\midrule
closure constants & $3$ / $3$ & $8.60$~s \\
array arguments & $1$ / $1$ & $2.67$~s \\
\bottomrule
\end{tabular}
\end{center}

The compile count is the load-independent quantity; the wall clock
matters at production scale, where a single f-block window compile is
of order fifteen minutes.  Three closures then cost about forty-five
minutes against fifteen minutes plus two evaluations in the array
form.

One restriction belongs to this section.  The geometry-dependent joint
entry \texttt{make\_nevpt2\_\allowbreak joint\_fn} cannot be traced under
\texttt{jax.jit} in this pyscfad build, because the SCF it runs inside
the call raises a tracer-conversion error.  The array-argument form is
therefore a fixed-geometry device and belongs to the theta-only
factories.  The correlation contraction itself is inside the jitted
theta-only function: the NEVPT2 energy in eager and jitted execution
agrees bit for bit, so the untraceable part is the SCF, not the
second-order correction.

\section{Repeated core folding in the CI matrix--vector product}

\texttt{core.ci\_matvec} builds the effective one-electron operator
once per matrix--vector product.  The original code recomputed that
fold separately inside the apply and diagonal paths, so every product
paid it twice.  The apply path now receives the folded quantities
$(h_{\mathrm{eff}}, g, e_{\mathrm{core}})$ from a shared helper, and
the diagonal path keeps its own call.  The output is unchanged bit for
bit: the norm of the test vector is $3977.4197449984595$ in every
interleaved run before and after the change, and the six leading
components agree as well.

The timing consequence is small.  In eager execution the apply
routine's cumulative time moved from $5.49$~s to $5.36$~s, about
$2.5$ percent and inside the load noise; after jitting, one product
takes $10$--$30$~ms and the before/after distributions are
indistinguishable.  This is a code-hygiene change, not a performance
promise.

\section{Choosing among the options}

\begin{center}
\begin{tabular}{p{0.30\textwidth}p{0.18\textwidth}p{0.42\textwidth}}
\toprule
Option & Default & When to change it \\
\midrule
\texttt{ci\_solver} & \texttt{'auto'} &
Keep the default.  Force \texttt{'plain'} only when the exact
pre-change CI vector must be reproduced; force \texttt{'blocked'} only
on a large determinant space where the plain tensors do not fit, and
expect the single-call time to grow. \\
MO transform & stepwise &
Always use \texttt{stepwise\_mo\_ints}; the one-shot form is not
exposed as an option. \\
jit reference arrays & array form &
Use \texttt{component\_args} and \texttt{arrays} for every repeated
geometry.  Do not jit the joint geometry-dependent entry. \\
\texttt{ci\_matvec} folding & single fold &
No user action; the change is internal and numerically neutral. \\
\bottomrule
\end{tabular}
\end{center}

Two properties of the measurements above belong to their reading.
Wall-clock numbers are relative and were taken on one shared node with
the load recorded; the compile counts, the peak resident memory, and
the energy and gradient differences are the reproducible quantities.
The 2000-determinant threshold is a choice between two measured
anchors, and a user whose determinant count lies between them should
measure the two routes on that space rather than assume the default.
